% ----------------------------------------------------------
% Introdução (exemplo de capítulo sem numeração, mas presente no Sumário)
% ----------------------------------------------------------
%\chapter*[Introdução]{Introdução}
\chapter{Contextualização}
%\addcontentsline{toc}{chapter}{INTRODUÇÃO}
% ----------------------------------------------------------

\section{Introdução}
O ensino experimental de Física em instituições públicas brasileiras enfrenta barreiras de escalabilidade devido ao alto custo de sistemas proprietários de aquisição de dados (DAQ). Embora o uso de smartphones pessoais tenha sido explorado, ele introduz falta de padronização e distrações no ambiente de aula. Este trabalho propõe uma arquitetura baseada em Microcontroladores (ESP32, Arduino) e Raspberry Pi que utiliza redes sem fio para centralizar a coleta de dados, permitindo que o docente gerencie múltiplas bancadas simultaneamente sem a necessidade de um computador dedicado por experimento.

\subsection{Contextualização do Tema}
A introdução deve começar apresentando o tema geral do seu trabalho. Aqui, você deve situar o leitor no contexto da sua pesquisa, mostrando a importância do tema em um cenário mais amplo. Use dados, informações atualizadas ou citações de autores para embasar o contexto.
\begin{itemize}
    \item Dica: Comece com uma frase impactante, uma estatística relevante ou um dado histórico que chame a atenção do leitor.
    \item Pergunta-chave: Qual é o tema geral do trabalho e por que ele é relevante?
\end{itemize}
\subsubsection*{Exemplo:}
“As mudanças climáticas têm sido consideradas um dos maiores desafios do século XXI, com impactos significativos sobre a saúde pública e os ecossistemas. No Brasil, esses impactos são ainda mais preocupantes devido à alta incidência de doenças transmitidas por vetores, como o Aedes aegypti.”

\subsection{Apresentação e Delimitação do Problema}
O problema central é a ociosidade e o isolamento do hardware em laboratórios didáticos. A pergunta de pesquisa foca na viabilidade técnica: Como construir uma infraestrutura de baixo custo e baseada em Wi-Fi que permita a centralização de dados experimentais em uma interface web única para fins pedagógicos?.

\section{Justificativa e Contribuições}
A literatura demonstra que laboratórios integrados aumentam a autonomia e o engajamento, mas a complexidade da infraestrutura ainda é um obstáculo. A contribuição deste trabalho é democratizar o acesso à instrumentação científica através de uma solução que reduz a dependência de cabos e minimiza a ociosidade de hardware.

\subsection{Pergunta de Pesquisa}
Formule a pergunta central que guia o seu estudo. A pergunta deve ser clara, objetiva e estar diretamente relacionada ao problema apresentado.
\begin{itemize}
    \item Dica: A pergunta deve ser respondida pelos objetivos do trabalho.
    \item Pergunta-chave: Qual é a questão principal que a pesquisa pretende responder?
\end{itemize}
\subsubsection*{Exemplo:}
“Como modelos de aprendizado de máquina podem ser aplicados para prever a distribuição do Aedes aegypti em regiões afetadas por mudanças climáticas?”

\section{Objetivos}

\subsection{Objetivo Geral}
\begin{quote}
Construir e apresentar um sistema de monitoramento centralizado para experimentos de Física, integrando microcontroladores a um servidor via rede Wi-Fi com interface gráfica.
\end{quote}
\section{Objetivos Específicos}

\begin{enumerate}
    \item Avaliar plataformas de hardware quanto ao custo para integração em experimentos didáticos de Física.
    \item Desenvolver a arquitetura de software (Firmware e Servidor Central) para telemetria via HTTP.
    \item Implementar uma Prova de Conceito (PoC) com o experimento de Carga e Descarga de Capacitor.
    \item Demonstrar a funcionalidade do sistema através da coleta de dados em tempo real.
\end{enumerate}
